% ============================================================
%  इस पुस्तक का उपयोग कैसे करें (How to use this book) + icon legend
% ============================================================
\chapter*{इस पुस्तक का उपयोग कैसे करें}
\markboth{इस पुस्तक का उपयोग कैसे करें}{इस पुस्तक का उपयोग कैसे करें}

यह पुस्तक \textbf{गतिविधि एवं परियोजना आधारित} है - इसका उद्देश्य AI को रटाना नहीं,
बल्कि \emph{स्वयं करके} सीखना है। हर अध्याय में सरल व्याख्या के साथ-साथ करने योग्य
कार्य, रोचक तथ्य एवं सोचने योग्य प्रश्न दिए गए हैं। नीचे पुस्तक में प्रयुक्त विभिन्न
बॉक्स (boxes) का अर्थ दिया गया है:

\vspace{8pt}
% रंगीन चिप + विवरण
\newcommand{\legendchip}[2]{\tikz[baseline=-0.55ex]\node[fill=#1,text=white,
  rounded corners=4pt,inner xsep=7pt,inner ysep=3pt,font=\cvHH\small]{#2};}

\begin{description}[leftmargin=0pt,labelsep=10pt,itemsep=9pt,
   font=\normalfont,style=nextline]
  \item[\legendchip{cvCyan}{गतिविधि}]
     कक्षा या घर पर स्वयं करने योग्य छोटा प्रयोग या कार्य।
  \item[\legendchip{cvBlue}{परियोजना}]
     समूह में करने योग्य बड़ा, बहु-चरणीय कार्य।
  \item[\legendchip{cvPurple}{क्या आप जानते हैं?}]
     विषय से जुड़ा रोचक एवं प्रेरक तथ्य।
  \item[\legendchip{cvOrange}{सोचिए}]
     चर्चा एवं चिंतन हेतु खुला प्रश्न।
  \item[\legendchip{aigreen}{मुख्य शब्द}]
     अध्याय के अंत में महत्वपूर्ण शब्दों का संक्षिप्त संग्रह।
  \item[\legendchip{cvOrange!85!red}{भारत में AI}]
     भारत में AI के वास्तविक उदाहरण एवं उपलब्धियाँ।
  \item[\legendchip{aiteal}{अग्रिम पठन}]
     विषय को और गहराई से समझने हेतु विश्वसनीय स्रोत।
\end{description}

\vspace{6pt}
\noindent\textbf{अन्य संकेत:}
\begin{itemize}[leftmargin=1.4em]
  \item प्रत्येक अध्याय के आरम्भ में \textbf{``इस अध्याय में आप सीखेंगे''} सूची एवं अंत में
        \textbf{सारांश}, \textbf{अभ्यास} तथा \textbf{बहुविकल्पीय प्रश्न (MCQ)} दिए गए हैं।
  \item वैज्ञानिक/तकनीकी शब्द \textbf{मोटे} अक्षरों में हिन्दी एवं कोष्ठक में
        अंग्रेज़ी \textit{(English)} दोनों रूप में दिए गए हैं; इनका अर्थ अंत की
        \textbf{शब्दावली} में भी है।
  \item \textbf{QR कोड} को मोबाइल कैमरे से स्कैन कर सीधे सम्बन्धित उपकरण, Colab notebook
        या प्रश्नोत्तरी तक पहुँचा जा सकता है।
  \item उत्तर पुस्तक के अंत में \textbf{उत्तर-कुंजी} में दिए गए हैं - पहले स्वयं प्रयास करें!
\end{itemize}

\vspace{4pt}
\noindent\emph{किसी भी गतिविधि हेतु महँगे उपकरण आवश्यक नहीं - एक मोबाइल/कंप्यूटर,
इंटरनेट और जिज्ञासा ही पर्याप्त है।}
